\documentclass{exam}
\usepackage{amsmath, amssymb}

\pagestyle{headandfoot}
\firstpageheadrule
\runningheadrule
\firstpageheader{Prof. Pham \\ Mathematical Theory of Probability}{Homework 8}{Aadhithya Saravanan}
\runningheader{Mathematical Theory of Probability \\ Homework 8}{}{Saravanan}
\firstpagefooter{}{}{}
\runningfooter{}{\thepage}{ }

\printanswers

\begin{document}

\underline{Problem 1}
\newline
\newline
Let $X$ represent the number of working components. The distribution of working parts on a rainy day is given by $\text{Bin}(n,p_1)$ and $\text{Bin}(n,p_2)$ for a dry day. We are looking for $P(X\ge k)$.
\[P(X\ge k)\]
\[=P(X\ge k|\text{Rain})P(\text{Rain})+P(X\ge k|\text{Dry})P(\text{Dry})\]
\[=P(\text{Bin}(n,p_1)\ge k)\alpha+P(\text{Bin}(n,p_2)\ge k)(1-\alpha)\]
\[=\alpha\sum_{i=k}^{n}\binom{n}{i}p_1^i(1-p_1)^{n-i}+(1-\alpha)\sum_{i=k}^{n}\binom{n}{i}p_2^i(1-p_2)^{n-i}\]
\newline

\underline{Problem 2}
\newline
\begin{parts}
    \part The probability of a correct decision is equal to the probability of acquitting an innocent defendant given the defendant is innocent plus the probability of convicting a guilty defendant given the defendant is guilty. The distribution of the number of votes for conviction for an innocent defendant is given by $\text{Bin}(12,0.1)$, which we want to be less than 9 to be correct. The distribution of the number of votes for conviction for a guilty defendant is given by $\text{Bin}(12,0.8)$, which we want to be at least 9 to be correct. So,
    \[P(\text{Correct})\]
    \[=P(\text{Correct|Innocent})P(Innocent)+P(\text{Correct|Guilty})P(Guilty)\]
    \[=0.35P(\text{Bin}(12,0.1)<9)+0.65P(\text{Bin}(12,0.8)\ge9)\]
    \[=0.35\sum_{k=0}^{8}\binom{12}{k}0.1^k0.9^{12-k}+0.65\sum_{k=9}^{12}\binom{12}{k}0.8^k0.2^{12-k}\]
    \part The probability of a conviction is equal to the probability of convicting an innocent defendant given the defendant is innocent plus the probability of convicting a guilty defendant given the defendant is guilty. This is the same as the last scenario, but we want at least 9 votes instead of less than 9 votes for the innocent defendant. So,
    \[P(\text{Conviction})\]
    \[=P(\text{Conviction|Innocent})P(Innocent)+P(\text{Conviction|Guilty})P(Guilty)\]
    \[=0.35P(\text{Bin}(12,0.1)\ge9)+0.65P(\text{Bin}(12,0.8)\ge9)\]
    \[=0.35\sum_{k=9}^{12}\binom{12}{k}0.1^k0.9^{12-k}+0.65\sum_{k=9}^{12}\binom{12}{k}0.8^k0.2^{12-k}\]
    \newline
\end{parts}

\underline{Problem 3}
\newline
\begin{parts}
    \part The chance of player 1 playing one game is equal to $\frac{2}{3}$, as player 1 has a $\frac{1}{3}$ chance of winning against player 2. To win two games, player 1 must win the first game and lose the second, giving a probability of $\frac{1}{3}\cdot\frac{3}{4}$. To win three games, player 1 must win the first game and second games, giving a probability of $\frac{1}{3}\cdot\frac{1}{4}$. Multiplying each of these probabilities by the corresponding number of games, we get
    \[E(X)\]
    \[=\frac{2}{3}\cdot1+\frac{1}{3}\cdot\frac{3}{4}\cdot2+\frac{1}{3}\cdot\frac{1}{4}\cdot3\]
    \[=\frac{17}{12}\]
    \part The chance of player 3 playing one game is equal to the probability of losing against the winner of the game between player 1 and player 2. If player 1 wins against player 2, which happens with probability $\frac{1}{3}$, player 3 has a $\frac{1}{4}$ chance of losing against player 1. If player 2 wins against player 1, which happens with probability $\frac{2}{3}$, player 3 has a $\frac{2}{5}$ chance of losing against player 2. This gives a probability of $\frac{1}{3}\cdot\frac{1}{4}+\frac{2}{3}\cdot\frac{2}{5}=\frac{7}{20}$. The probability of player 3 winning two games is equal to the probability of winning the first game, which is $1-\frac{7}{20}=\frac{13}{20}$. Multiplying each of these probabilities by the corresponding number of games, we get
    \[E(X)\]
    \[=\frac{7}{20}\cdot1+\frac{13}{20}\cdot2\]
    \[=\frac{33}{20}\]
    \newline
\end{parts}

\underline{Problem 4}
\newline
\newline
The probability of no error in the article is equal to the probability of an error by the first typist given that it is typed by the first typist plus the probability of an error by the second typist given that it is typed by the second typist. The probability of an error by a typist is modeled by a Poisson distribution with a $\lambda$ equal to the average number of errors for the typist. So,
\[P(\text{No Error})\]
\[=P(\text{No Error}|\text{First Typist})P(\text{First Typist})+P(\text{No Error}|\text{Second Typist})P(\text{Second Typist})\]
\[=0.5P(\text{Poisson}(3)=0)+0.5P(\text{Poisson}(4.2)=0)\]
\[=0.5(e^{-3}\cdot\frac{3^0}{0!})+0.5(e^{-4.2}\cdot\frac{4.2^0}{0!})\]
\[=0.5(e^{-3}+e^{-4.2})\]
\newline

\underline{Problem 5}
\newline
\begin{parts}
    \part Let $X$ represent the number of people hit by lightning in a month. It is modeled by $\text{Poisson}(\frac{400000}{100000})=\text{Poisson}(4)$. We want $P(X\ge8)$. So,
    \[P(X\ge8)\]
    \[=1-P(X<8)\]
    \[=1-\sum_{k=0}^{7}e^{-4}\cdot\frac{4^k}{k!}\]
    \part Let $Y$ represent the number of months with 8 or more people being hit by lightning. $Y$ follows a binomial distribution with $n=12$ and $p=1-\sum_{k=0}^{7}e^{-4}\cdot\frac{4^k}{k!}$. We want $P(Y\ge2)$. So,
    \[P(Y\ge2)\]
    \[=1-P(Y<2)\]
    \[=1-\sum_{k=0}^{k=1}\binom{12}{k}p^k(1-p)^{12-k}\]
    with $p=1-\sum_{k=0}^{7}e^{-4}\cdot\frac{4^k}{k!}$.
    \part Let $Z$ represent the distribution of the first month that 8 more people being hit by lightning. It follows a geometric distribution with $p=1-\sum_{k=0}^{7}e^{-4}\cdot\frac{4^k}{k!}$. We want $P(Z=i)$ for $i\ge1$. So,
    \[P(Z=i)\]
    \[=(1-p)^{i-1}p\]
    with $p=1-\sum_{k=0}^{7}e^{-4}\cdot\frac{4^k}{k!}$. This assumes that each month is independent of the others and each month has the same distribution of lightning hitting people.
    \newline
\end{parts}

\underline{Problem 6}
\newline
\newline
This follows a negative binomial model, with the winning team needing 3 wins in the first $i-1$ games. So, we are looking for $P(X=i)$, so
\[P(X=i)\]
\[=\binom{i-1}{3}0.6^3(0.4)^{i-3}\]
Summing these, we get a probability of winning of 0.710.
\newline
\newline
The probability of winning in a 3-game series is $\binom{3}{2}0.6^2(0.4)+0.6^3=0.648<0.710$.
\newline

\underline{Problem 7}
\newline
\begin{parts}
    \part With 5 people, the problem follows a binomial model with $n=5$ and $p=\frac{2}{3}$. We are looking for $P(X\ge5)=(\frac{2}{3})^5$.
    \part With 8 people, the problem follows a binomial model with $n=8$ and $p=\frac{2}{3}$. We are looking for $P(X\ge5)=\sum_{k=5}^{8}(\frac{2}{3})^k(1-\frac{2}{3})^{8-k}$.
    \part This follows a negative binomial mode $Y$ with $n=5$ and $p=\frac{2}{3}$. We are looking for $P(Y=6)$, which is equal to $\binom{5}{4}p^5(1-p)$.
    \part This is the same as the last part with $n=6$. We are looking for $P(Y=7)$, which is equal to $\binom{6}{4}p^5(1-p)^2$ with $p=\frac{2}{3}$.
    \newline
\end{parts}

\end{document}